% Propietario conservado: /Users/ruben/Documents/New project/output/EXTREMA_MEDIA_RAZON_RECIPROCIDAD_ES_EN_20260918/ARTICULO/ES/source/sections/04_accion_realizacion.tex
% Procedencia: integral 2249, §25.10--13 y cap. 89; propietarios 02_pantalla_accion,
% 37_cuerdas_branas_teoria_m, terna_espectral_teoria_m, capitulo_25_pantallas_teoria_m.
\section{Action sur les trajectoires et réalisation dynamique undécadimensionnelle}
\label{x_reciprocidad:iv:accion-realizacion}

L’extension de la structure à une dynamique dimensionnelle requiert de
transporter conjointement l’état, les canaux et les données de
frontière. Cette section commence par une fonctionnelle définie sur
les trajectoires du TPK et par sa loi exacte de composition.
La surface d’univers et le codomaine undécadimensionnel sont
introduits ensuite, avec les applications qui doivent conserver cette
composition et les opérateurs déjà construits.

\subsection{Descente d’une évolution et suffisance de la donnée observée}

Soit \(q:\mathcal X_{\mathrm{HMT}}\to Y\) une lecture
surjective des états avec mémoire, et soit
\(U:\mathcal X_{\mathrm{HMT}}\to\mathcal X_{\mathrm{HMT}}\)
une évolution complète. Une évolution sur les valeurs
observées existe précisément lorsque la donnée conservée par
\(q\) suffit à déterminer la valeur observée suivante.

\begin{theorem}[Critère exact de descente par fibres]
\label{x_reciprocidad:iv:thm:descenso-fibras}
Il existe une unique application \(\overline U:Y\to Y\)
satisfaisant \(qU=\overline Uq\) si et seulement si
\begin{equation}
 q(x)=q(y)\quad\Longrightarrow\quad q(Ux)=q(Uy).
 \label{x_reciprocidad:iv:eq:descenso-fibras}
\end{equation}
\end{theorem}
\begin{proof}
Si \(\overline U\) existe, appliquer la même fonction à
\(q(x)=q(y)\) prouve la condition. Réciproquement, on
définit \(\overline U(q(x))=q(Ux)\). La condition
\eqref{x_reciprocidad:iv:eq:descenso-fibras} garantit l’indépendance du
représentant. La surjectivité de \(q\) définit
l’application sur tout \(Y\) et prouve son unicité.
\end{proof}

Le théorème permet de décider quelle mémoire une
lecture doit conserver. Lorsque deux représentants ont des images observées
différentes, on peut élargir la donnée retenue, restreindre
le domaine ou utiliser une correspondance multivoque.
Ce sont des choix mathématiques différents. En particulier, le
retour d’une coordonnée de phase n’autorise pas à identifier
les évolutions de deux histoires qui conservent des mémoires
différentes.

\subsection{Cocycle des changements directionnels et arithmétiques}

Une trajectoire admissible du TPK porte, sur chaque arête
\(e_j\), une direction \(d_j\) et un type arithmétique
\(t_j\). Pour compter les changements dans une concaténation,
on conserve également la donnée précédente. L’arête augmentée est
\(\widetilde e_j=(e_j,d_{j-1},t_{j-1})\) ; la paire
\((d_0,t_0)\) fait partie de la frontière initiale. Nous définissons
\[
 w_{\mathrm{giro}}(\widetilde e_j)
       =\mathbf1_{\{d_j\ne d_{j-1}\}},\qquad
 w_{\mathrm{tipo}}(\widetilde e_j)
       =\mathbf1_{\{t_j\ne t_{j-1}\}}.
\]
Pour \(\widetilde\gamma=(\widetilde e_1,\ldots,
\widetilde e_R)\), soient
\begin{equation}
 \mathbf I(\widetilde\gamma)=
 \left(\sum_{j=1}^R w_{\mathrm{giro}}(\widetilde e_j),
       \sum_{j=1}^R w_{\mathrm{tipo}}(\widetilde e_j)\right)
 =\bigl(N_{\mathrm{giro}},N_{\mathrm{tipo}}\bigr).
 \label{x_reciprocidad:iv:eq:cociclo-accion}
\end{equation}
Deux trajectoires se composent lorsque leurs extrémités coïncident
et que la donnée précédente de la seconde coïncide avec la donnée
terminale de la première.

\begin{proposition}[Additivité avec frontière conservée]
\label{x_reciprocidad:iv:prop:accion-aditiva}
Pour toute paire composable,
\begin{equation}
 \mathbf I(\widetilde\gamma\widetilde\eta)
 =\mathbf I(\widetilde\gamma)+\mathbf I(\widetilde\eta).
 \label{x_reciprocidad:iv:eq:accion-aditiva}
\end{equation}
Par conséquent, \(\mathbf I\) est un cocycle additif à valeurs
dans \(\mathbb N^2\) sur la catégorie des trajectoires
augmentées. Pour chaque \(\lambda>0\), la fonctionnelle
\[
 I_\lambda(\widetilde\gamma)
 =N_{\mathrm{giro}}(\widetilde\gamma)
      +\lambda N_{\mathrm{tipo}}(\widetilde\gamma)
\]
est aussi additif.
\end{proposition}
\begin{proof}
La concaténation ne change la donnée précédente d’aucune
arête : celle de la première arête de \(\widetilde\eta\)
coïncide, par la condition de composition, avec la donnée terminale
de \(\widetilde\gamma\). La somme sur toutes les arêtes
se divise exactement en les deux sommes de
\eqref{x_reciprocidad:iv:eq:accion-aditiva}. Appliquer la fonctionnelle linéaire
\((a,b)\mapsto a+\lambda b\) prouve la dernière affirmation.
\end{proof}

L’information de frontière fait effectivement partie de la preuve.
Si deux segments d’une arête sont évalués séparément
en ignorant leur raccord, tous deux peuvent avoir un compte
de tournants nul même si leurs directions sont différentes. Le
chemin concaténé contient alors un tournant. Conserver
\((d_{j-1},t_{j-1})\) enregistre ce terme dans la première
arête du second segment et restitue l’additivité.

Pour des extrémités et des données de frontière fixées, le principe
d’interaction minimale du corpus sélectionne des minimiseurs de
\(I_\lambda\) parmi les routes admissibles. Le principe
est une règle de sélection, tandis que le résultat
suivant établit une propriété d’existence finie.

\begin{proposition}[Existence finie de minimiseurs]
Si l’ensemble \(\Gamma(a,b)\) des trajectoires admissibles
à bord fixé est fini et non vide, il existe
\(\gamma_*\in\Gamma(a,b)\) tel que
\(I_\lambda(\gamma_*)\leq I_\lambda(\gamma)\)
pour tout \(\gamma\in\Gamma(a,b)\).
\end{proposition}
\begin{proof}
L’image de \(\Gamma(a,b)\) par \(I_\lambda\) est
un ensemble fini non vide de nombres réels, qui possède
un élément minimum et au moins une préimage.
\end{proof}

La proposition ne sélectionne pas la valeur de \(\lambda\)
et ne confond pas l’existence avec l’unicité. La valeur et la règle
de sélection utilisées dans une réalisation déterminée
doivent faire partie de ses données déclarées. La fonctionnelle
additive retrouvée ici quantifie les changements du chemin
lui-même ; son transport vers une action dimensionnelle constitue
une opération ultérieure.

\subsection{Persistance étendue et surface d’univers}

La réalisation d’une route comme objet étendu distribue
sa généalogie sur une coordonnée interne. Pour une corde,
l’histoire possède deux coordonnées locales \(\sigma^a\) :
l’une spatiale interne et l’autre temporelle. Une réalisation se
décrit par une application \(X:\Sigma\to\mathcal M\), une
métrique \(h_{ab}\) sur la surface d’univers et des champs
de fond dans l’espace cible. Les conditions sur
\(\partial\Sigma\) font partie du domaine variationnel.
Pour une brane de dimension spatiale \(p\), l’histoire
est un volume d’univers de dimension \(p+1\).

Le codomaine variationnel de corde utilisé par le corpus
est l’action de Polyakov
\begin{equation}
 S_{\mathrm P}[X,h]
 =-\frac{T_{\mathrm s}}2\int_\Sigma d^2\sigma\,
 \sqrt{-h}\,h^{ab}g_{MN}(X)
             \partial_aX^M\partial_bX^N,
 \label{x_reciprocidad:iv:eq:polyakov}
\end{equation}
où \(T_{\mathrm s}\) est la tension. La notation
\(T_{\mathrm s}\) la distingue de la dualité
\(\mathsf T_0\), du transport cardinal \(T_d\)
et de l’évolution complète \(U\).

La composition des trajectoires augmentées possède déjà la
propriété démontrée dans \ref{x_reciprocidad:iv:prop:accion-aditiva}.
Une réalisation continue de cette composition doit
spécifier l’application des chemins vers les configurations,
son transport des conditions de frontière et la
normalisation qui compare \(I_\lambda\) à l’action
continue. L’intégrale \eqref{x_reciprocidad:iv:eq:polyakov} précise
le codomaine et ses champs ; ses paramètres ne choisissent pas
rétroactivement les émissions ni les registres du TPK.

\subsection{Réalisation à onze dimensions et opérateurs gradués}

La réduction dirigée par \(K\) fournit
\(V_{11}=\ell_K\oplus V_{10}\), et le théorème
\ref{x_reciprocidad:iv:thm:pantallas-isomorfas} coordonne cette réduction
avec l’écran réticulaire et la dualité. La réalisation
physique est typée au moyen de
\[
 \Xi_M:\mathscr S_M^{\mathrm{alg}}
                  \longrightarrow\mathcal X_{11}^{\mathrm{phys}}.
\]
Le codomaine porte une métrique \(g_{MN}\), une forme de degré trois
\(C_{MNP}\), un gravitino \(\psi_M\), une action,
des contraintes de jauge et des transformations de
supersymétrie. Une réalisation opératorielle inclut
\[
 \mathfrak S_M=(\mathcal H_{\mathrm{phys}},H,Q),\qquad
 \mathcal H_{\mathrm{phys}}=\mathcal H_+\oplus\mathcal H_-.
\]
La graduation s’exprime au moyen d’une involution
\(\Gamma\), égale à \(+I\) sur \(\mathcal H_+\)
et à \(-I\) sur \(\mathcal H_-\). La supercharge
est impaire, \(\Gamma Q=-Q\Gamma\), et
l’hamiltonien est positif ou nul. La relation
\(Q^2=H\) s’entend comme une égalité d’opérateurs,
y compris leurs domaines.

Soit \(\mathcal H_{\mathrm{HMT}}\) une réalisation
hilbertienne des états généalogiques pertinents,
avec des opérateurs \(\widehat H,\widehat Q\). Une
isométrie \(W:\mathcal H_{\mathrm{HMT}}\to
\mathcal H_{\mathrm{phys}}\) doit transporter
les domaines, la graduation et les observables. Les conditions
opératorielles s’écrivent
\begin{equation}
 W\widehat H=HW,\qquad
 W\widehat Q=QW,\qquad
 WU_T^{\mathrm{HMT}}=U_T^{\mathrm{phys}}W.
 \label{x_reciprocidad:iv:eq:entrelazamiento-fisico}
\end{equation}

\begin{theorem}[Transport compatible de la supercharge]
\label{x_reciprocidad:iv:thm:supercarga-transporte}
Supposons \(\widehat Q^{\,2}=\widehat H\),
\(W D(\widehat Q)\subset D(Q)\), et que les
entrelacements de \eqref{x_reciprocidad:iv:eq:entrelazamiento-fisico}
sont satisfaits dans leurs domaines. Alors, pour
\(\psi\in D(\widehat H)=D(\widehat Q^{\,2})\),
\[
 W\psi\in D(Q^2),\qquad Q^2W\psi=HW\psi.
\]
Si, en outre, le transport physique des rayons respecte
\eqref{x_reciprocidad:iv:eq:conjugacion-h}, le secteur compact
transporté conserve son spectre sous des rayons réciproques.
\end{theorem}
\begin{proof}
Comme \(\psi\in D(\widehat Q^2)\), aussi bien
\(\psi\) que \(\widehat Q\psi\) appartiennent
à \(D(\widehat Q)\). Par conséquent, \(W\psi\in D(Q)\)
et \(QW\psi=W\widehat Q\psi\in D(Q)\). On peut
appliquer \(Q\) à nouveau, et
\[
 Q^2W\psi=W\widehat Q^{\,2}\psi
          =W\widehat H\psi=HW\psi.
\]
L’égalité spectrale du secteur compact est transportée
par l’isométrie et l’entrelacement unitaire de T ;
sur l’image fermée de \(W\), celui-ci est une
équivalence unitaire des réalisations correspondantes.
\end{proof}

Le théorème conserve explicitement les opérateurs et le
domaine sur lequel l’identité est prouvée. La
graduation organise deux secteurs ; la construction des
champs fermioniques et de leurs relations d’occupation possède
sa propre structure. Les distributions thermiques et le
catalogue des espèces ne sont utilisés ni dans la démonstration
des écrans ni dans la dualité compacte de cet
article.

\subsection{Compactification compatible et conclusion du développement}

Une compactification de l’état se décrit au moyen
d’une section prolongeable \(x_\infty\), d’une application
\(\kappa_{\mathrm{comp}}\) vers l’espace de modules
choisi et d’un transport \(\mathcal U\) des données
de métrique, d’orientation, de frontière et de mémoire. L’indice
distingue cette application de la lecture rationnelle
\(\kappa\) du registre dodécaphasique. Sa compatibilité
exige quatre propriétés concrètes :
\begin{enumerate}
\item les restrictions finies de \(x_\infty\)
      appartiennent aux prolongements admissibles du TPK ;
\item l’application de compactification entrelace les dualités,
      \(\kappa_{\mathrm{comp}}\mathsf T
        =\mathsf T_M\kappa_{\mathrm{comp}}\) ;
\item les observables et le spectre descendent aux
      quotients choisis conformément au critère
      \ref{x_reciprocidad:iv:thm:descenso-fibras} ;
\item la réalisation transporte la direction
      \(\ell_K\), les conditions de frontière et les
      domaines opératoires, et conserve les équations
      d’action et de supersymétrie déclarées.
\end{enumerate}

Les preuves précédentes fournissent la réalisation
hilbertienne de phase, la fermeture traciale, les écrans
avec leurs applications de quotient, l’involution de rayon,
son équivalence unitaire et le cocycle additif des
trajectoires. La réalisation undécadimensionnelle
spécifie les champs et les compatibilités qu’elle
utilise lorsqu’elle reçoit cette structure. Son identification
physique complète conserve en outre les conditions
de tension, d’oscillateurs, d’amplitudes, d’anomalies et de
limite de basse énergie énumérées par le corpus.
Ces conditions n’altèrent pas les identités
algébriques démontrées : elles délimitent l’affirmation
ultérieure obtenue en les transportant.

Le prolongement exceptionnel et le prolongement dimensionnel
restent liés par leur origine commune. Le
Monstre agit dans la construction graduée
\(V^\natural\) ; le Hamiltonien compact et les
écrans reçoivent les données d’autres applications du
même état. Cette formulation conserve le
contenu positif de la double réalisation,
avec ses opérations, ses preuves et l’information
transmise à chaque étape (théorème~\ref{x_reciprocidad:iv:thm:supercarga-transporte}).
